%% ============================================================
%% abstract.tex  —  Thesis Abstract
%% ============================================================
%% Drop-in note: this file provides the abstract text inside a
%% standard `abstract` environment.  If your thesis class instead
%% expects an unnumbered chapter, replace the environment with:
%%   \chapter*{Abstract}
%%   \addcontentsline{toc}{chapter}{Abstract}
%%   ... text ...
%% ============================================================

\chapter*{Abstract}
\addcontentsline{toc}{chapter}{Abstract}

Saturn's moon \titan{} is the only known world beyond Earth with stable
surface liquids, a dense organic-rich atmosphere, and a subsurface water
reservoir, making it a site where prebiotic chemistry can be studied at
planetary scale over geological time.  Previous habitability assessments
have been largely qualitative or restricted to a single location or epoch.
This thesis develops the first quantitative, spatially-resolved, multi-epoch
habitability map of \titan{}'s surface, combining eight Cassini observational
proxy features (liquid hydrocarbon, organic abundance, chemical energy,
methane cycling, surface--atmosphere exchange, topographic complexity,
geomorphological diversity, and subsurface water access) within a Bayesian
Beta-conjugate framework.  The posterior habitability probability
$P(H \mid \featurevec)$ is evaluated at every pixel of a global grid
(\SI{4490}{\metre\per\text{px}}, 6.49 million pixels) at five temporal anchor
epochs from the Late Heavy Bombardment ($-3.5\gya$) to the red-giant ocean
($+5.9\gya$), with 95\% credible intervals at each location.

At the present epoch, north polar lake margins are the highest-scoring
environments: Towada Lacus ranks first ($P(H) = 0.573$), driven by the highest
spectrally resolved organic abundance of any polar lake combined with surface
liquid and active methane cycling.  The temporal trajectory is non-monotonic,
passing through three regimes: organic-accumulation dominance during the
bombardment, polar-lake dominance from $-1.0\gya$ through the present, and a
global water--ammonia ocean from $+5.1\gya$ in which former polar basins again
lead (Uvs Lacus, $P(H) = 0.713$).  A transient solvent-free interval of
$\sim$1.1\,Gyr separates the hydrocarbon and water-ocean regimes.  Applied as
an external check, the framework independently identifies Selk crater, the
\dragonfly{} landing site, as the most geomorphologically diverse equatorial
location ($f_7 = 0.629$, $6.8\times$ the global median) from Cassini data
alone.  Throughout, the maps quantify the co-occurrence of the physical and
chemical preconditions for life, not the presence of life or the likelihood of
its origin.

\vspace{1em}
\noindent\textbf{Keywords:} Cassini; astrobiology; ocean worlds; prebiotic chemistry; methane lakes; subsurface ocean;
Selk crater; VIMS; synthetic aperture radar; planetary science; acetylenotrophy; \titan{}; Beta-conjugate; 
habitability; Bayesian inference; multi-epoch; temporal evolution; \dragonfly{}
